% Models i mètodes d'interpretabilitat
\section{Funcionament Tobii}\label{sec:tobii}

El hardware de Tobii (Spectrum) enregistra quina zona de la pantalla miren els ulls a 600Hz.
El seu software calcula les fixacions, i permet definir cada paraula d'un text com a àrea d'interés (AOI en les seves sigles en anglès).

El problema a l'hora d'exportar les dades és que Tobii Pro Lab no optimitza el format.
En comptes de crear un fitxer per a cada text, crea un fitxer amb tots els temps com a línia i totes les AOI com a columna.
Això fa que el tamany del fitxer creixi quadràticament amb el nombre de textos i el programa es bloquegi.
Per tal d'optimitzar aquest procés hem exportat diversos fitxers de la següent manera\footnote{Dins el menú Analyze/Data export}.
Utilitzem sempre el filtre \textit{Tobii I-VT (Fixation)} i \textit{noise reduction} tant per l'obertura de l'ull com per diàmetre de la pupi"la.

\begin{enumerate}
    \item Dades generals, inclou dia, hora i resultats de la calibració. (General, excepte AOI size, timestamp i event value; un sol Time of interest; cap AOI)
    \item Anotacions. Els inputs per teclat dels participants. (Participant name, Event, Event value; Times of Interest = Media gaze data; cap AOI)
    \item Dades. Cal seleccionar grups petits d'AOI i els seus respectius temps a Times of interest. Finalment es van fer 4 grups de 10 textos per cada participant. Es poden exportar de cop a fitxers individuals amb \textit{Format=multiple standard files}.
        \begin{itemize}
            \item General: Participant name
            \item Eye tracking data: Pupil diameter, Validity of eye data
            \item Media: stimulus name
            \item Eye movement events: eye movement type, duration and index, fixation point, AOI hit
        \end{itemize}
\end{enumerate}

\section{Models de llenguatge}

En aquest treball utilitzem diferents tipus de models basats en l'arquitectura \textit{transformer} 
\cite{vaswaniAttentionAllYou2017}, que ha revolucionat el processament del llenguatge natural (PLN). 

\subsection{Models BERT finetunejats}

Els models BERT (\textit{Bidirectional Encoder Representations from Transformers}) són models 
\textit{encoder-only} que han demostrat un rendiment excepcional en tasques de classificació de text.
En el nostre cas, utilitzem variants en castellà com \textit{PlanTL-GOB-ES/roberta-large-bne} i 
el més recent \textit{BSC-LT/MrBERT-es}.

El procés de \textit{finetuning} consisteix a adaptar aquests models preentrenats a la nostra tasca 
específica de detecció de sexisme. Vegeu el codi per a més detalls sobre la implementació, que 
inclou la preparació de les dades, la configuració del model i l'entrenament utilitzant la 
llibreria \textit{transformers} de HuggingFace.

\todo{revisar termes IA termcat}
\subsection{Models de llenguatge extensos}

A més dels models finetunejats, també experimentem amb grans models de llenguatge (LLMs) 
en configuracions de \textit{zero-shot} i \textit{few-shot}. Aquests enfocaments permeten 
avaluar la capacitat dels models per generalitzar sense necessitat d'un entrenament 
específic per a la nostra tasca.

En \textit{zero-shot}, el model realitza la classificació sense cap exemple 
previ de la tasca. En \textit{few-shot}, es proporcionen uns quants exemples (normalment 
entre 1 i 5) per guiar la predicció del model. Vegeu el codi per a la implementació concreta 
d'aquests enfocaments.

\section{Mètodes d'interpretabilitat}

Per comprendre els patrons d'atenció dels models i comparar-los amb els patrons humans, 
utilitzem la llibreria Captum \cite{captum}, que proporciona una suite completa de mètodes d'atribució.

\subsection{Mètodes funcionals}

Els mètodes que hem implementat són:

\begin{itemize}
    \item \textbf{Saliency}: Mètode basat en gradients que calcula la derivada de la sortida 
    respecte als \textit{embeddings} d'entrada. És útil per identificar quins tokens tenen 
    major impacte en la decisió del model.
    
    \item \textbf{Integrated Gradients}: Aquest mètode calcula la integral dels gradients 
    al llarg d'un camí des d'una línia base (normalment zeros) fins a la entrada real. 
    Proporciona una atribució més robusta que el \textit{saliency} simple i inclou una 
    mètrica de convergència (\textit{delta}) que indica la qualitat de l'atribució.
    
    \item \textbf{Layer Grad-CAM}: Combina informació de gradients amb mapes d'activació 
    de capes intermèdies del model. Tot i que la nostra implementació actual retorna 
    principalment zeros, aquest mètode té potencial per identificar quines capes i 
    neurones són més rellevants per a la decisió.
\end{itemize}

\subsection{Mètodes en desenvolupament}

Els següents mètodes estan implementats però encara no funcionen correctament a causa de 
limitacions tècniques:

\begin{itemize}
    \item \textbf{DeepLift}: Mètode que compara l'activació de cada neurona amb una 
    referència, proporcionant una atribució basada en la diferència respecte a aquesta 
    línia base.
    
    \item \textbf{Guided Backpropagation} i \textbf{Guided Grad-CAM}: Mètodes que 
    combinen \textit{backpropagation} guiat amb mapes d'activació per obtenir 
    visualitzacions més interpretables.
    
    \item \textbf{Input X Gradient}: Multiplica l'entrada pels gradients per obtenir 
    una mesura de la importància de cada característica.
    
    \item \textbf{Noise Tunnel}: Afegeix soroll a les entrades i mitja els resultats 
    de múltiples execucions per obtenir atribucions més robustes.
\end{itemize}

\subsection{Mètodes massa lents}

\begin{itemize}
    \item \textbf{Feature Ablation}: Aquest mètode elimina característiques i mesura 
    l'impacte en la sortida, però és computacionalment costós a causa de la 
    alta dimensionalitat dels \textit{embeddings} (768 dimensions).
\end{itemize}

\subsection{Implementació pràctica}

El procés d'interpretabilitat segueix les següents etapes (vegeu el codi per a la 
implementació detallada):

\begin{enumerate}
    \item Carregar el model entrenat i el tokenitzador corresponent
    \item Tokenitzar el text d'entrada i obtenir els \textit{embeddings}
    \item Aplicar el mètode d'atribució seleccionat (p. ex., Integrated Gradients)
    \item Sumar les atribucions al llarg de la dimensió dels \textit{embeddings} per obtenir 
    una puntuació per a cada token
    \item Visualitzar els resultats utilitzant codis de colors (verd per a contribucions 
    positives, vermell per a negatives)
\end{enumerate}

El codi inclou funcions per generar visualitzacions HTML amb els tokens ressaltats segons 
la seva importància per a la predicció del model.
\todo{afegir exemples}

\section{Comparació amb patrons humans}

L'objectiu final d'aquest treball és comparar els patrons d'atenció obtinguts amb aquests mètodes amb els patrons d'atenció 
humans registrats amb l'\textit{eye-tracker}. Aquesta comparació permetrà:

\begin{itemize}
    \item Avaluar si els models se centren en les mateixes parts del text que els humans
    \item Identificar possibles biaixos o limitacions en els models computacionals
    \item Explorar si les divergències entre atenció humana i computacional es relacionen 
    amb errors de classificació
    \item Obtenir \textit{insights} sobre com millorar els models per alinear-los millor amb 
    els processos cognitius humans
\end{itemize}

Aquesta anàlisi comparativa es durà a terme utilitzant les dades d'\textit{eye-tracking} 
recollides amb el dispositiu Tobii, tal com es descriu a la Secció\ref{sec:tobii}, 
i els resultats s'analitzaran quantitativament i qualitativament per obtenir 
conclusions sobre la similitud entre ambdós tipus d'atenció.

\section{Disponibilitat del codi}

Tot el codi utilitzat en aquest treball, incloent-hi la implementació dels models, els 
mètodes d'interpretabilitat i els scripts d'anàlisi, està disponible públicament 
al repositori GitHub: \url{https://github.com/Pastells/eye\_tracker\_sexism}. 
